\documentclass[../../book]{subfiles}

\begin{document}

\textbf{Motivation.} Cryptography is wonderful. If you are holding this book, you are probably 
already know the significance of cryptography in the modern 
informational infrastructure: from classical standard protocols such as 
Transport Layer Security (TLS) to the more recent cryptographic
advancements used in Blockchain technologies. This book is 
exactly about the latter.

In particular, as the book's name suggests, we consider the \emph{zero-knowledge
cryptography}, which is a cornerstone of numerous privacy-preserving protocols.
However, in its current state, the zero-knowledge technology is
\emph{tremendously} hard to work with:
\begin{itemize}
    \item The technology is relatively new, so there are almost 
    no high-level abstractions that can be used to simplify the
    design of zero-knowledge protocols for regular developers.
    \item The technology is based on advanced mathematics, which is hard to
    understand without a proper background.
    \item The amount of available resources on the topic is surprisingly limited.
    What's more, the resources are very diverse: they are either (a) too
    high-level, which makes you understand what the protocol \emph{should} do,
    but you still have no idea how to implement it, (b) or too low-level,
    consisting of numerous security proofs and obscure theorems with little
    practical implications. Note that in the latter case you typically still
    have no clue what was going on, unless you spend a significant amount of
    time studying the topic.
\end{itemize}

Of course, this book cannot solve all the aforementioned problems at once.
However, we do hope that this book will serve as a complete, practical guide to
most state-of-the-art techniques in zero-knowledge cryptography for practicing
engineers. We gathered all the necessary information in one place, and tried to
make it easy-to-follow, with a lot of examples and code snippets. Yet, we still
try to preserve the mathematical rigor where suitable and necessary.

\textbf{Who are we?} We are \href{https://blockstream.com/}{\emph{Blockstream}}
--- a global leader in Bitcoin and blockchain technologies. 
Our research and engineering efforts were focused on numerous topics, ranging 
from confidential transactions~\cite{FC:PBFMW18} to threshold- and multisignatures~\cite{C:NicRufSeu21}.
Additionally, the majority of authors at the time of writing this book were employed
at \href{https://distributedlab.com/}{\emph{Distributed Lab}}, which developed 
zero-knowledge solutions involving optimizing advanced algebraic
system components and social applications with the focus on privacy (e.g.
\href{https://rarimo.com/}{\emph{Rarimo}}).

\textbf{Who is this book for?} This book does require basic background in
Mathematics. Of course, we will try to explain all the necessary concepts from
scratch (such as basic number and group theory, security analysis), but this
material is rather supplementary and is not the main focus of the book. That
being said, you do not need Math PhD, but be prepared: the book is challenging.

\textbf{How to use this book?} The book is structured as follows:
\begin{enumerate}
    \item The first part of the book, consisting of \Cref{section:number-theory} to \Cref{section:commitment-schemes}, covers the mathematical and cryptographic background needed to understand zero-knowledge protocols.
    If you feel comfortable with polynomials, groups, elliptic curves, and commitment schemes, you can skip this part.
    If not, this is a perfect place to get a basic understanding of these topics.
    Keep in mind that they are useful well beyond the scope of this book.
    \item The main part of the book starts with \Cref{section:intro-zk}, where we define zero-knowledge proofs and explain the basic concepts behind them.
    Then we consider $\Sigma$-protocols (\Cref{section:sigma}), the simplest yet one of the most fundamental families of zero-knowledge protocols.
    \item We then proceed to more advanced topics: zk-SNARKs (\Cref{section:r1cs} to \Cref{section:circom}), PlonK (\Cref{section:plonk}), and Bulletproofs (\Cref{section:bulletproofs}).
    \item From this point on, the book becomes progressively more challenging.
    We first consider the sum-check protocol and its applications (\Cref{section:sumcheck,section:gkr}), together with lookup arguments (\Cref{section:toolkit,section:lookups,section:ultragroth}).
    Then we move to post-quantum schemes, both lattice-based (\Cref{section:lattices,section:lattice-sigma-protocols,section:labrador}) and code-based, including the famous ZK-STARK (\Cref{section:linear-codes,section:iopp,section:zk-stark}).
\end{enumerate}

\end{document}
